% --- Chapter: Binary Expressions and Operators ---
\section{Binary Expressions and Operators}

CRISP provides a rich set of binary operators organized by precedence. The parser uses a Pratt parsing algorithm with defined binding powers for each operator, eliminating ambiguity in complex expressions.

\subsection{Operator Precedence Table}

Operators are evaluated from highest to lowest binding power. Higher in the table = tighter binding:

\begin{table}[H]
\centering
\rowcolors{2}{softviolet!30}{white}
\caption{\textcolor{darkpurple}{Operator Precedence — Highest to Lowest}}
\vspace{0.3em}
\begin{tabular}{@{}cllp{5cm}@{}}
\toprule
\rowcolor{softvioletdark}
\textbf{\textcolor{white}{Precedence}} & \textbf{\textcolor{white}{Operators}} & \textbf{\textcolor{white}{Associativity}} & \textbf{\textcolor{white}{Description}} \\
\midrule
1 (highest) & \texttt{**} & Right & Exponentiation \\
2 & \texttt{* / \%} & Left & Multiplication, division, modulo \\
3 & \texttt{+ -} & Left & Addition, subtraction, string concatenation \\
4 & \texttt{..} & Non-associative & Range construction \\
5 & \texttt{< > <= >= <=>} & Non-associative & Comparison, spaceship \\
6 & \texttt{== != =~ !~} & Non-associative & Equality, regex match \\
7 & \texttt{\&\&} & Left & Logical AND \\
8 & \texttt{||} & Left & Logical OR \\
9 (lowest) & \texttt{=} & Right & Assignment \\
\bottomrule
\end{tabular}
\end{table}

\subsection{Arithmetic Operators}

\begin{tcolorbox}[colback=codebg, colframe=darkpurple, colbacktitle=darkpurple, title=\textbf{\textcolor{codegold}{Arithmetic Operators}}, fonttitle=\bfseries, sharp corners]
\begin{lstlisting}
# Basic arithmetic
let sum = 10 + 5;        # 15
let diff = 10 - 5;       # 5
let prod = 10 * 5;       # 50
let quot = 10 / 5;       # 2
let mod = 10 % 3;        # 1

# Exponentiation (right-associative)
let pow = 2 ** 3;        # 8
let chain = 2 ** 3 ** 2; # 512 (2 ** (3 ** 2) = 2 ** 9)

# Mixed types: Int + Float → Float
let mixed = 10 + 3.5;    # 13.5 (Float)

# Division always produces Float
let half = 10 / 4;       # 2.5 (Float, not truncated)

# Negative numbers
let neg = -5 + 3;        # -2
\end{lstlisting}
\end{tcolorbox}

\subsection{String Concatenation}

The \texttt{+} operator concatenates strings. At least one operand must be a string; the other is converted via \texttt{.as\_str()}:

\begin{tcolorbox}[colback=codebg, colframe=darkpurple, colbacktitle=darkpurple, title=\textbf{\textcolor{codegold}{String Concatenation}}, fonttitle=\bfseries, sharp corners]
\begin{lstlisting}
# String + String
let greeting = "Hello, " + "World!";
say greeting;  # Hello, World!

# String + Number
let msg = "Value: " + 42;
say msg;  # Value: 42

# String + Bool
say "Flag: " + true;  # Flag: true

# Number + String works too
say 100 + " points";  # 100 points

# Chaining
let full = "CRISP" + " v" + 0 + "." + 1 + "." + 0;
say full;  # CRISP v0.1.0
\end{lstlisting}
\end{tcolorbox}

\subsection{Comparison Operators}

Comparison operators return \texttt{true} or \texttt{false}. They are non-associative — chaining them without parentheses is an error:

\begin{tcolorbox}[colback=codebg, colframe=darkpurple, colbacktitle=darkpurple, title=\textbf{\textcolor{codegold}{Comparison Operators}}, fonttitle=\bfseries, sharp corners]
\begin{lstlisting}
# Numeric comparisons
say 10 > 5;        # true
say 10 < 5;        # false
say 10 >= 10;      # true
say 10 <= 5;       # false

# Equality (works across types)
say 10 == 10;      # true
say 10 == 5;       # false
say 10 != 5;       # true
say "hi" == "hi";  # true

# Floats compare as f64
say 3.14 == 3.14;  # true

# Non-associative — use && for chaining
# say 1 < x < 10;     # Error!
say 1 < x && x < 10;  # Correct

# Null equality
let x = null;
say x == null;     # true
say null != 42;    # true
\end{lstlisting}
\end{tcolorbox}

\subsection{Spaceship Operator}

The \texttt{<=>} (spaceship) operator returns \texttt{-1}, \texttt{0}, or \texttt{1} for less-than, equal, or greater-than:

\begin{tcolorbox}[colback=codebg, colframe=darkpurple, colbacktitle=darkpurple, title=\textbf{\textcolor{codegold}{Spaceship Operator}}, fonttitle=\bfseries, sharp corners]
\begin{lstlisting}
# Numeric comparison
say 3 <=> 5;    # -1  (less)
say 5 <=> 5;    # 0   (equal)
say 7 <=> 5;    # 1   (greater)

# String comparison (lexicographic)
say "apple" <=> "banana";  # -1
say "crisp" <=> "crisp";   # 0
say "zebra" <=> "apple";   # 1

# Useful for sorting callbacks
let sorted = sort(arr, |a, b| => a <=> b);
\end{lstlisting}
\end{tcolorbox}

\subsection{Logical Operators}

\texttt{\&\&} (AND) short-circuits — if the left operand is falsy, the right is never evaluated. \texttt{||} (OR) short-circuits — if the left operand is truthy, the right is never evaluated:

\begin{tcolorbox}[colback=codebg, colframe=darkpurple, colbacktitle=darkpurple, title=\textbf{\textcolor{codegold}{Logical Operators}}, fonttitle=\bfseries, sharp corners]
\begin{lstlisting}
# AND (&&) — both must be true
if x > 0 && x < 100 {
    say "x is in range";
}

# OR (||) — at least one must be true
if name == "admin" || name == "root" {
    say "privileged user";
}

# Short-circuit: expensive() is never called
let x = 0;
if x != 0 && expensive() {
    say "won't happen";
}

# Short-circuit: fallback() is never called
let y = 42;
if y > 0 || fallback() {
    say "short-circuited";
}

# Combined complex conditions
if (age >= 18 && age <= 65) || has_permission {
    say "access granted";
}

# Truthiness: false, null, 0, "" are falsy
if 0 || null || false || "" {
    say "never runs";
} else {
    say "all falsy";  # this runs
}
\end{lstlisting}
\end{tcolorbox}

\subsection{Range Operator}

The \texttt{..} operator creates a range. Ranges are used primarily in \texttt{for} loops and array slicing:

\begin{tcolorbox}[colback=codebg, colframe=darkpurple, colbacktitle=darkpurple, title=\textbf{\textcolor{codegold}{Range Operator}}, fonttitle=\bfseries, sharp corners]
\begin{lstlisting}
# Numeric range (inclusive start, exclusive end)
for i in 0..5 {
    say i;
}
# Output: 0 1 2 3 4

# Range stored in variable
let r = 1..10;

# Array slicing with range
let @arr = [10, 20, 30, 40, 50];
say @arr[1..3];  # [20, 30]

# String slicing
say "CRISP"[0..3];  # "CRI"
\end{lstlisting}
\end{tcolorbox}

\subsection{Regex Operators}

CRISP provides \texttt{=~} (match) and \texttt{!~} (negated match) for regular expression testing:

\begin{tcolorbox}[colback=codebg, colframe=darkpurple, colbacktitle=darkpurple, title=\textbf{\textcolor{codegold}{Regex Operators}}, fonttitle=\bfseries, sharp corners]
\begin{lstlisting}
let text = "The quick brown fox";

# =~ returns boolean in scalar context
if text =~ /fox/ {
    say "Found fox!";
}

# !~ negated match
if text !~ /dragon/ {
    say "No dragons here";
}

# With g flag, =~ returns array of matches
let words = text =~ /\w+/g;
say words;  # ["The", "quick", "brown", "fox"]

# Case-insensitive
if text =~ /FOX/i {
    say "case-insensitive match";
}

# Combine with logical operators
if text =~ /quick/ && text =~ /fox/ {
    say "Contains both quick and fox";
}
\end{lstlisting}
\end{tcolorbox}

\subsection{Assignment Operator}

The \texttt{=} operator reassigns an existing variable. It is right-associative, enabling chained assignment:

\begin{tcolorbox}[colback=codebg, colframe=darkpurple, colbacktitle=darkpurple, title=\textbf{\textcolor{codegold}{Assignment Operator}}, fonttitle=\bfseries, sharp corners]
\begin{lstlisting}
# Reassignment (variable must already exist)
let x = 10;
x = 20;       # valid — x exists
x = x + 1;    # valid — self-referencing

# Right-associative chained assignment
let a = 0;
let b = 0;
a = b = 42;
say a;  # 42
say b;  # 42

# Assignment in expression context
let y = (x = 5) + 10;
say x;  # 5
say y;  # 15

# Error: assigning to undeclared variable
# z = 5;  # Error: Undefined variable: z
\end{lstlisting}
\end{tcolorbox}

\subsection{Ternary Operator}

The ternary \texttt{? :} operator provides inline conditional evaluation. It is right-associative and has the lowest precedence above assignment:

\begin{tcolorbox}[colback=codebg, colframe=darkpurple, colbacktitle=darkpurple, title=\textbf{\textcolor{codegold}{Ternary Operator}}, fonttitle=\bfseries, sharp corners]
\begin{lstlisting}
# Basic ternary
let label = score > 50 ? "pass" : "fail";

# Ternary with expressions
let max = a > b ? a : b;

# Nested ternary (right-associative)
let grade = score >= 90 ? "A"
          : score >= 80 ? "B"
          : score >= 70 ? "C"
          : "F";

# Ternary in expressions
say "You " + (age >= 18 ? "may" : "may not") + " enter.";
\end{lstlisting}
\end{tcolorbox}

\begin{tcolorbox}[colback=lightlavender, colframe=softvioletdark, title=\textbf{\textcolor{darkpurple}{Design Note: Pratt Parsing}}, fonttitle=\bfseries, sharp corners]
CRISP uses a Pratt (top-down operator precedence) parser for binary expressions. Each operator is assigned a left and right binding power:

\begin{itemize}
    \item \textbf{Left binding power:} Determines whether this operator "grabs" the expression on its left from the previous operator. If the previous operator's right power is lower than this operator's left power, this operator wins.
    \item \textbf{Right binding power:} Determines whether this operator "allows" the next operator to grab its right operand. Higher right power = tighter binding on the right.
    \item \textbf{Associativity:} Implemented by making left binding power lower than right (left-associative, e.g., \texttt{a - b - c} = \texttt{(a - b) - c}) or equal (right-associative, e.g., \texttt{a = b = c} = \texttt{a = (b = c)}).
\end{itemize}

This approach eliminates the need for explicit precedence tables in the grammar — precedence is encoded directly in the numeric binding powers.
\end{tcolorbox}
